% !TEX program = lualatex
\documentclass[aspectratio=43,professionalfonts,handout]{beamer}
\usepackage{beamer_template}

\newcommand{\LectureNum}{10}
\newcommand{\LectureTitle}{鋭角の三角比の応用}
\newcommand{\TextPages}{p.\,128-129}
\newcommand{\NextTopic}{鈍角の三角比}
\newcommand{\NextPages}{p.\,129-132}
\newcommand{\CourseName}{基礎数学B}
\renewcommand{\TermName}{後期}

\begin{document}

% -----------------------------------------------
% スライド1: 表紙＋目標＋用語・定義
% -----------------------------------------------
\begin{frame}{\CourseName\quad 第\LectureNum 回「\LectureTitle 」}
  {\small \TermName\quad \TeacherName\quad ／\quad 教科書 \TextPages\quad
  \textbf{目標}：三角比を用いて図形の辺の長さを求められる}

  \begin{block}{用語・定義}
  \begin{description}[labelwidth=5.5em]
    \item[\kw{仰角・俯角}] 仰角：水平面から上方向への角度（見上げる角）；俯角：水平面から下方向への角度（見下ろす角）
    \item[\kw{三角比の応用（距離・高さを求める）}] 直角三角形を作図して $\sin$, $\cos$, $\tan$ のどれを使うか判断する
  \end{description}
  \end{block}
\end{frame}

% -----------------------------------------------
% スライド2: 公式・性質
% -----------------------------------------------
\begin{frame}{公式・性質}
  \begin{block}{}
  \begin{itemize}
    \item 三角比の活用手順：図を描く→関係式を立てる→計算
    \item 三角関数表の使い方
  \end{itemize}
  \end{block}
\end{frame}

% -----------------------------------------------
% スライド3: 例1→問1
% -----------------------------------------------
\begin{frame}{例1→問1}
  \begin{exampleblock}{例題1) p.130（高さ・距離の計算）}
    （板書で解説）
  \end{exampleblock}

  \vfill

  \begin{block}{問3) p.130\quad 自分で解いてみよう}
    （ノートに解くこと）
  \end{block}
\end{frame}

% -----------------------------------------------
% スライド4: 例2→問2
% -----------------------------------------------
\begin{frame}{例2→問2}
  \begin{exampleblock}{例題2) p.131（2点間の距離）}
    （板書で解説）
  \end{exampleblock}

  \vfill

  \begin{block}{問4) p.131\quad 自分で解いてみよう}
    （ノートに解くこと）
  \end{block}
\end{frame}

% -----------------------------------------------
% まとめ
% -----------------------------------------------
\begin{frame}{まとめ}
  \begin{itemize}
    \item 仰角・俯角を用いて高さや距離を三角比で表す
    \item 直角三角形を作図して $\sin$, $\cos$, $\tan$ のどれを使うか判断する
    \item 答えは $\sqrt{}$ を有理化して整理する
  \end{itemize}

  \vfill
  {\small 基本問題：272, 273, 274}
\end{frame}

\end{document}
